% 第5回　リサーチ・クエスチョンの確定（記入例・模範解答）
\session{5}{1}{10月28日（水）}{リサーチ・クエスチョンの確定}{AIによる批判的検証}{}

\goals{
  \item AIを「厳しい査読者」として使い、問いの弱点を洗い出せる
  \item 反論に答える形で、問いの背景と意義を論理的に説明できる
  \item 前半の成果を、AI活用記録を添えたレポートにまとめられる
}

\work{準備}{AIに提示する材料}
\abox[仮の問い]{14mm}{名古屋都市圏の20代にとって、オンラインでの情報収集は、自動車ディーラーへの来店のしかたをどう変えたのか。}
\abox[背景（200字程度）]{28mm}{若い世代の「車離れ」が言われる一方で、同世代は車の情報を動画や比較サイトで詳しく調べている。先行研究では、購入前の情報探索が多い人ほど訪問する販売店が少ないことが示されている（第4回の文献）。しかし、車の利用が多い地方都市圏での状況や、来店の目的の変化はまだ十分に調べられていない。販売店にとっても、来店の意味の変化をつかむことは店舗のあり方を考える手がかりになる。}
\abox[主な先行研究（3件）]{20mm}{① 第4回で精読した文献（購入前の情報探索と訪問店舗数）　② 日本自動車工業会「乗用車市場動向調査」　③ 総務省「通信利用動向調査」}

\work[70mm]{ワーク1・2}{査読者の指摘に、補強か修正で答える}
\begin{promptbox}[査読者として批判させる]
あなたは厳しい査読者です。次のリサーチ・クエスチョンについて、1. 半年で調べて答えられるか　2. 先行研究に対する新しさがあるか　3. 範囲は適切か　4. 用語は明確か　の観点から問題点を指摘し、それぞれ改善案を示してください。問い：（あなたの問い）　背景：（200字程度）　主な先行研究：（3件）
\end{promptbox}
\begin{wstabh}{colspec={Q[l,wd=22mm,m] X[3,l,m] Q[c,wd=14mm,m] X[3,l,m]},row{2-Z}={ht=17mm}}
  観点 & AIの指摘（要点） & 補強／\newline 修正 & 自分の対応と、その根拠 \\
  答えられるか & \af{「変えたのか」は過去との比較が必要で、半年では難しい} & \ans{修正} & \af{過去との比較はやめ、「来店前に何を調べ、何のために来店するか」の現状を調べる問いにする} \\
  新しさ & \af{先行研究と同じことを調べるだけにならないか} & \ans{補強} & \af{先行研究は「訪問店舗数」。自分は「来店の目的」と「車の利用が多い地域」に注目する点が新しい} \\
  範囲 & \af{20代全体は広い。購入予定のない人も含むのか} & \ans{修正} & \af{「この1年に車の購入・買い替えを検討した20代」に絞る} \\
  用語の明確さ & \af{「来店のしかた」が曖昧} & \ans{修正} & \af{「来店の回数」と「来店の目的（見る・試乗・相談・契約）」に分けて定義する} \\
\end{wstabh}

\begin{promptbox}[反論に答える練習]
私の問いに対して「（AIが挙げた反論）」という批判があります。この批判に答える論点を3つ挙げ、それぞれどんな根拠が必要かを示してください。根拠が見つかっていない論点は、そのように明示してください。
\end{promptbox}

\work{修正の記録}{問いはどう変わったか}
\atwoboxes{修正前の問い}{修正後の問い（確定）}{24mm}
{名古屋都市圏の20代にとって、オンラインでの情報収集は、自動車ディーラーへの来店のしかたをどう変えたのか。}
{車の購入を検討した名古屋都市圏の20代は、来店前にオンラインで何を調べ、何を目的にディーラーへ来店しているのか。}

\work[90mm]{ワーク3}{確定した問いを1枚にまとめる（発表用）}
\begin{wstab}{colspec={Q[l,wd=30mm,m] X[l,m]},row{1}={ht=14mm},row{2-Z}={ht=22mm}}
  \hc{問い} & \af{車の購入を検討した名古屋都市圏の20代は、来店前にオンラインで何を調べ、何を目的にディーラーへ来店しているのか。} \\
  \hc{背景と意義} & \af{若い世代は車の情報をオンラインで詳しく調べるようになった。来店の意味が「比べる場」から変わっているなら、販売店の役割や店づくりも変わる必要がある。車の利用が多い地域の若者の実態を知ることに意義がある。} \\
  \hc{先行研究\newline でわかって\newline いること} & \af{都市部では、情報探索が多い人ほど訪問店舗数は少なく、1回の商談は長い（第4回の文献）。公的調査から、若い世代のインターネット利用の広がりがわかる。来店の目的と地域差は未解決。} \\
  \hc{調べ方の見通し} & \af{① 公的調査で全体の傾向を確認する　② 購入を検討した20代の知人・学生に、来店前の情報収集と来店の目的をアンケートで聞く（20人程度）　③ 可能ならディーラーの担当者に話を聞く} \\
\end{wstab}

\work{発表}{1人3分。受けた質問と答え}
\atwoboxes{受けた質問}{答え・今後の課題}{30mm}
{20人のアンケートで「20代」と言えるのか？\newline 「来店の目的」の選択肢はどう決めるのか？}
{一般化はできないので、「傾向を探る調査」と位置づけ、限界として書く。選択肢は先行研究の分類を参考にし、自由記述欄も設ける。}

\work[50mm]{提出物（前半）}{リサーチ・クエスチョン確定レポートの要件チェック}
\begin{achecks}
  \item 確定したリサーチ・クエスチョンと、選んだ理由
  \item 背景と意義
  \item 先行研究の整理（実在を確認した文献リスト付き）
  \item 調べ方の見通し
  \item AI活用記録（AI利用ガイドの形式）
  \item 本文の冒頭または末尾に、AI利用を明記する一文
\end{achecks}
\afillin{書式・締切（授業内で案内）}{（授業内で案内した書式・締切を記入）}
\aimemoA{問いへの批判（査読者役）と、反論に答える論点の整理}
{「あなたは厳しい査読者です。4つの観点から問題点と改善案を」「この批判に答える論点を3つ、必要な根拠とともに」}
{4つの指摘のうち3つで問いを修正。新しさの指摘には、先行研究との違いを示して反論した。改善案の文言はそのまま使わず、自分で書き直した。}
{修正後の問いが、手元の文献とアンケートで答えられるかを、先行研究の調査項目と照らして確かめた}

\begin{point}
  \item 問いの修正の記録（修正前と修正後）が書けているかが最も大切。AIの指摘を受けて、何をどう変えたかがレポートの「AI活用記録」の中心になる。
  \item すべての指摘を受け入れる必要はない。新しさの指摘のように、根拠を示して反論できた点は高く評価する。
  \item 「〜をどう変えたのか」のような過去との比較を含む問いは、学部2年の半年では答えにくい。現状を記述する問いに絞り込めているかを見る。
  \item 発表では、調べ方の見通しと限界（対象の少なさなど）を自分で言えるかを見る。
\end{point}
